Este capítulo presenta un análisis crítico e integrador de los resultados obtenidos durante el desarrollo y evaluación del sistema de clasificación histopatológica prostática basado en aprendizaje profundo y \textit{Multiple Instance Learning} (MIL). La discusión se centra en la interpretación del comportamiento de las distintas arquitecturas evaluadas, sus implicaciones clínicas, los aportes en términos de interpretabilidad y las principales limitaciones metodológicas y computacionales identificadas.

\section{Síntesis e Interpretación de los Hallazgos Principales}
Los resultados experimentales ponen de manifiesto diferencias claras y consistentes entre las estrategias MIL evaluadas. En particular, las arquitecturas basadas en mecanismos de atención (ABMIL y SmABMIL) superaron de forma sistemática a los enfoques de agregación no adaptativa (MeanMIL y MaxMIL), tanto en métricas discriminativas globales como en estabilidad frente a distintas particiones del conjunto de datos.\newline

Desde una perspectiva metodológica, este comportamiento sugiere que la capacidad de aprender pesos de relevancia a nivel de instancia constituye un factor determinante para abordar la heterogeneidad espacial inherente a las imágenes histopatológicas. A diferencia de los métodos de pooling clásico, los modelos de atención permiten concentrar el proceso de decisión en subconjuntos informativos del bag, mitigando el efecto de regiones no representativas o ruidosas.\newline

En el ámbito clínico, SmABMIL destacó por ofrecer un equilibrio favorable entre sensibilidad, especificidad y AUC, acompañado de una menor variabilidad inter-fold. Este patrón sugiere una mayor capacidad de generalización frente a la diversidad morfológica presente entre pacientes, un aspecto crítico en aplicaciones diagnósticas reales.\newline

Por el contrario, MaxMIL evidenció un sesgo pronunciado hacia la clase negativa, reflejado en una elevada tasa de falsos negativos. Este comportamiento limita su aplicabilidad clínica, dado que la omisión de lesiones malignas clínicamente significativas representa un riesgo inaceptable en escenarios de cribado y diagnóstico asistido.

\section{Análisis Comparativo de Estrategias de Agregación MIL}
El análisis comparativo entre las distintas arquitecturas permite identificar limitaciones estructurales inherentes a los esquemas de agregación no adaptativa:

\begin{itemize}
    \item \textbf{MeanMIL} tiende a diluir señales discriminativas provenientes de regiones tumorales pequeñas al promediar de forma uniforme todas las instancias del bag.
    \item \textbf{MaxMIL} depende de manera excesiva de una única instancia extrema, amplificando la influencia de ruido, artefactos o regiones no representativas.
\end{itemize}

En contraste, los modelos basados en atención introducen un mecanismo de selección suave que permite asignar pesos diferenciados a las instancias, adaptando la representación del bag a la distribución espacial de patrones histopatológicos relevantes. Esta propiedad se refleja directamente en una mejora del desempeño clínico y en una mayor estabilidad durante el entrenamiento.\newline

La variante SmABMIL incorpora, además, un mecanismo de atención compuerta que actúa como regularizador implícito, evitando distribuciones de atención excesivamente concentradas y favoreciendo representaciones más robustas y clínicamente coherentes.

\section{Implicaciones Clínicas del Enfoque Propuesto}
Desde una perspectiva diagnóstica, los resultados obtenidos respaldan el uso de modelos MIL basados en atención para el análisis de biopsias prostáticas digitalizadas, donde la señal patológica suele estar localizada en regiones específicas y no necesariamente dominantes en términos de área.\newline

En este contexto, la priorización de métricas como la sensibilidad y el AUC resulta particularmente relevante, dado que la detección temprana de lesiones clínicamente significativas tiene un impacto directo en la toma de decisiones terapéuticas. El comportamiento observado en SmABMIL sugiere un compromiso adecuado entre la detección de casos positivos y el control de falsos positivos, alineándose con las necesidades clínicas reales.\newline

Adicionalmente, la disponibilidad de mecanismos de interpretabilidad basados en atención posiciona al sistema propuesto como un potencial apoyo a la decisión clínica, permitiendo al patólogo contrastar visualmente la predicción del modelo con regiones tisulares específicas de interés.

\section{Interpretabilidad y Correlato Histopatológico}
Uno de los aportes metodológicos más relevantes de este trabajo es la integración explícita de análisis de interpretabilidad a nivel de instancia y de WSI. La inspección de los pesos de atención, tanto mediante histogramas como a través de mapas de calor espaciales, evidenció que los modelos ABMIL y SmABMIL concentran su atención en regiones coherentes con patrones histopatológicos asociados a malignidad.\newline

La proyección de los pesos de atención sobre las coordenadas originales de la WSI permitió generar mapas visuales que facilitan el correlato entre la decisión algorítmica y la estructura tisular subyacente. Asimismo, el análisis de los parches Top-K más relevantes mostró la capacidad del modelo para identificar regiones con alteraciones glandulares y patrones celulares anómalos.\newline

Estos resultados refuerzan la transparencia del enfoque propuesto y lo alinean con los principios de explicabilidad exigidos en aplicaciones clínicas basadas en inteligencia artificial.

\section{Limitaciones del Estudio y Líneas de Mejora}
A pesar de los resultados obtenidos, el presente estudio presenta una serie de limitaciones que delimitan el alcance de sus conclusiones y abren oportunidades para trabajos futuros.

\subsection{Gestión de Imágenes Gigapíxel y Escalabilidad}
El manejo de imágenes WSI, incluso tras su fragmentación en parches, impuso restricciones significativas en términos de almacenamiento, transferencia y procesamiento. La extracción de características y el entrenamiento de modelos MIL con bolsas de tamaño variable requirieron estrategias específicas de procesamiento por lotes para garantizar la viabilidad computacional del sistema.

\subsection{Variabilidad Cromática y Heterogeneidad del Tejido}
Las variaciones en la tinción H\&E introducen una fuente de ruido que puede afectar la robustez del aprendizaje. Aunque se aplicaron técnicas de filtrado y normalización básicas, este fenómeno sigue representando un desafío abierto en patología digital y podría abordarse mediante métodos avanzados de normalización de tinción o aprendizaje auto-supervisado.

\subsection{Restricciones Computacionales y Arquitecturales}
El entrenamiento de arquitecturas basadas en atención con bolsas de gran tamaño impuso limitaciones de memoria en GPU. Si bien se implementaron optimizaciones para mitigar estos efectos, la escalabilidad del enfoque sigue condicionada por la dimensionalidad de las características y el número de instancias por bag.

\subsection{Consideraciones Metodológicas}
Aunque se adoptaron particiones estrictas a nivel de paciente para evitar fuga de datos, el tamaño limitado del dataset y la formulación binaria del problema representan simplificaciones del escenario clínico real. Estudios futuros podrían explorar esquemas multiclase, análisis longitudinales o validaciones externas para fortalecer la validez clínica del sistema.\newline

Aunque estas limitaciones no invalidan los hallazgos presentados, sí establecen un marco claro para la interpretación de los resultados y para la evolución futura del enfoque propuesto.
